Learning-augmented (``with-predictions'') algorithms for online
bipartite matching, where requests arrive one at a time and each must be
matched to a resource or dropped before the rest of the input is seen,
have proliferated. One family consumes a prediction of how contended
each resource will be. Another tests a prediction of the arrival mix on
a short prefix of the requests --- a vanishing fraction of the instance
--- before committing to follow it or to fall back on an advice-free
baseline. These algorithms have been analyzed in isolation. We give the
first unified experimental study. We place all of them on one harness: a
common structured prediction-error model, a common optimum, and
confidence intervals. We run it on synthetic graphs, six real-world
graphs, and real request traces. Our central finding is that on
average-case inputs the value of predictions is robustness insurance,
not a performance lever. The advice-free baseline already reaches
\(0.99\) of the offline optimum on our few-types instances, leaving
under \(0.01\) for good advice to add, whereas unguarded
prediction-following crashes to \(0.47\) on the same instances; the
practical worth of the sophisticated algorithms is that they never
crash. We close with a theoretical outlook --- one proved
consistency/robustness trade-off inequality, and a quantitative reading
of our own experiments under which the prefix needed to decide whether
to trust advice scales as the inverse square of the advice's upside. The
full formal development is separate work, and no theorem beyond the
stated inequality is claimed here. Our experiments deliver one message
and the outlook explains why it should be expected: on average-case
matching, the advice's upside is smaller than the price of finding out
whether to trust it.
